% ----------------------------------------------------------
% Capítulo 3 — Metodologia
% ----------------------------------------------------------
\chapter{Metodologia}
\label{cap:metodologia}

Este capítulo descreve o método de trabalho adotado para o aprimoramento da solução de busca em linguagem natural sobre o acervo cartográfico da DSG. O trabalho parte de um protótipo de referência desenvolvido pelo 1º CGEO, descrito e analisado em detalhe no Capítulo~\ref{cap:aprimoramento}, e adota a estratégia de \textbf{substituição dirigida}: identificam-se os componentes responsáveis pelas limitações documentadas pelo cliente e substituem-se esses componentes por equivalentes mais robustos, preservando o que já funciona (esquema do banco, conjunto de parâmetros de busca, lógica das consultas espaciais).

O capítulo apresenta os elementos transversais ao projeto: o escopo (entregáveis e exclusões), os modelos comparados, a construção e a auditoria do \textit{dataset} de avaliação, os dois lotes de validação, as especificações v2 e v3 e o protocolo de medição. A execução do aprimoramento --- análise do protótipo, requisitos, planejamento da migração, desenvolvimento por camadas e cronograma --- é apresentada no Capítulo~\ref{cap:aprimoramento}.

\paragraph{Foco da avaliação: a camada de tradução linguística.} O objeto de avaliação deste trabalho é o desempenho do LLM na \textbf{tradução de uma consulta em linguagem natural para os parâmetros estruturados} de busca no catálogo. A execução dessa consulta contra o banco (PostgreSQL/PostGIS), a resolução de nomes geográficos e a interpretação das \textit{bounding boxes} das cartas são componentes herdados do protótipo do 1º CGEO; sua correção é assumida como \textit{caixa-preta} e não é medida. Métricas sobre o \textit{conjunto de produtos retornados} ficam fora do escopo (Seção~\ref{sec:exclusoes-avaliacao}) e são apontadas como direção de trabalho futuro no Capítulo~\ref{cap:conclusao}. O LLM é comparado contra um \textit{ground truth} de \textbf{parâmetros extraídos}, não de produtos recuperados.

% ---
\section{Escopo do Projeto}
% ---

\subsection{Entregáveis}

O projeto entregou:

\begin{itemize}
  \item uma nova arquitetura de \textit{backend} em Python (FastAPI + LangChain), que substitui o \textit{backend} Node.js do protótipo e implementa \textit{Tool Calling} nativo via Ollama;
  \item um \textit{dataset} de avaliação com 310 consultas de referência (\textit{ground truth}), distribuídas em seis categorias de consulta e uma de consultas fora do domínio, acompanhado de um manual de anotação e de um relatório de auditoria;
  \item um \textit{pipeline} de avaliação automatizado, que executa o \textit{dataset} sobre cada modelo e produz as métricas de acurácia, precisão, \textit{recall}, F1-\textit{score} e latência, com os testes estatísticos correspondentes;
  \item um relatório comparativo entre quatro modelos de LLM com suporte a \textit{Tool Calling}, executados localmente, com a recomendação fundamentada do modelo mais adequado para uso operacional sobre o acervo da DSG, complementado por uma referência paralela com modelos maiores servidos em nuvem;
  \item dois lotes de validação independentes (Seções~\ref{sec:lote-validacao} e~\ref{sec:v3}) e a v3 do \textit{Tool Calling}, com ferramentas auxiliares e validação com retorno, desenvolvida sobre as 310 consultas e o primeiro lote e testada no segundo contra um controle com Saída Estruturada; a v3 está implementada no pacote e no avaliador (\texttt{src/pfc\_busca/v3.py}), mas ainda não integrada à API do \textit{backend}.
\end{itemize}

\subsection{Exclusões}
\label{sec:exclusoes-avaliacao}

Estão fora do escopo deste trabalho:

\begin{itemize}
  \item \textbf{Avaliação da qualidade de recuperação em nível de produto}: uma vez extraídos os parâmetros, o conjunto de produtos retornado depende do motor de busca (\texttt{buscar\_catalogo}) e da qualidade do banco cartográfico, ambos herdados do protótipo. O trabalho \textbf{não mede} precisão e \textit{recall} em nível de produto. A execução SQL permanece no \textit{pipeline} apenas como demonstração de que a arquitetura fecha ponta a ponta;
  \item \textbf{Resolução avançada de nomes geográficos}: o cruzamento entre \textit{strings} livres (estado, município) e as \textit{bounding boxes} das cartas depende de recursos externos (malha do IBGE, Nominatim/OSM) e da qualidade das geometrias no banco; é tratado como trabalho futuro (Capítulo~\ref{cap:conclusao});
  \item desenvolvimento de interface gráfica (\textit{frontend}) para o usuário final;
  \item implantação em produção em ambiente da DSG ou de suas OM subordinadas;
  \item pesquisa de opinião ou avaliação qualitativa com usuários;
  \item treinamento ou ajuste fino (\textit{fine-tuning}) de modelos de linguagem;
  \item busca por coordenadas geográficas brutas (\textit{bounding box}) sem referência nominal a um lugar;
  \item suporte a idiomas além do português brasileiro.
\end{itemize}

\subsection{Requisitos que norteiam o escopo}
\label{sec:requisitos-resumo}

O escopo é norteado por requisitos funcionais (RF) e não funcionais (RNF) elicitados a partir da análise do protótipo e das conversas com o cliente. O Quadro~\ref{quadro:requisitos-resumo} traz o resumo; a discussão completa está no Capítulo~\ref{cap:aprimoramento} (Seção~\ref{sec:requisitos}).

\begin{quadro}[h]
\centering
\caption{Resumo dos requisitos que norteiam o escopo}
\label{quadro:requisitos-resumo}
\small
\begin{tabular}{|c|p{5.5cm}|c|p{5.0cm}|}
\hline
\multicolumn{2}{|l|}{\textbf{Requisitos Funcionais (RF)}} & \multicolumn{2}{l|}{\textbf{Requisitos N\~ao Funcionais (RNF)}} \\
\hline
RF1 & Receber consulta em linguagem natural via REST & RNF1 & Execu\c c\~ao totalmente local (privacidade) \\
\hline
RF2 & Traduzir consulta em par\^ametros do \textit{schema} & RNF2 & Lat\^encia mediana operacional aceit\'avel \\
\hline
RF3 & Executar busca no PostgreSQL/PostGIS & RNF3 & Conformidade a ET-PCDG e Perfil MGB \\
\hline
RF4 & Selecionar modelo por requisi\c c\~ao (\textit{dev only}) & RNF4 & Reprodutibilidade dos experimentos \\
\hline
RF5 & Registrar tra\c cos para as m\'etricas & RNF5 & Manutenibilidade (m\'odulos, testes) \\
\hline
    &                                       & RNF6 & Substituibilidade de modelo por configura\c c\~ao \\
\hline
\end{tabular}%
\fonte{Elaborado pelos autores. Detalhes no Cap.~\ref{cap:aprimoramento}, Seção~\ref{sec:requisitos}.}
\end{quadro}

Requisitos ideais são mensuráveis, mas alguns dos RNF --- em particular o RNF2, latência --- dependem de fatores externos ao escopo deste PFC (\textit{hardware} do ambiente de implantação, carga concorrente). Por isso os valores de latência são reportados nas condições de cada ambiente de execução (Seção~\ref{sec:ambiente}), sem limiar operacional fixado \textit{a priori}.

% ---
\section{Seleção de Modelos}
\label{sec:selecao-modelos}
% ---

\subsection{Critérios de Seleção}

Os modelos foram selecionados segundo os critérios: (i) disponibilidade no Ollama com suporte nativo a \textit{Tool Calling}, verificada no catálogo oficial da biblioteca~\cite{ollama2023}, que mantém rankings de popularidade e uma etiqueta explícita de compatibilidade com \textit{tool calling}; (ii) disponibilidade em versões quantizadas compatíveis com estações de trabalho modestas; (iii) lançamento ou atualização entre 2024 e 2026; e (iv) diversidade de arquiteturas e organizações desenvolvedoras.

\subsection{Modelos Avaliados}

Foram avaliados quatro modelos abertos, executados localmente pelo Ollama:

\begin{enumerate}
  \item \textbf{Qwen 3 4B}~\cite{qwen2025qwen3} (Alibaba Cloud), \textit{tag} \texttt{qwen3:4b-instruct-\allowbreak 2507-q4\_K\_M}: 4 bilhões de parâmetros, quantização Q4\_K\_M, 2,3~GB --- a menor \textit{build} dos quatro; mesmo assim, carregado com o contexto, não coube inteiro nos 4~GB de VRAM da estação de referência, onde o Ollama alocou na GPU \res{estacao}{qwen}{nagpu} do modelo carregado (Tabela~\ref{tab:ambiente_estacao}). Usou-se a \textit{build} oficial sem raciocínio interno (\textit{Instruct-2507}), pela razão explicada na Seção~\ref{sec:config-modelos};
  \item \textbf{Gemma 4 E4B}~\cite{team2026gemma4} (Google DeepMind), \textit{tag} \texttt{gemma4:\allowbreak e4b-it-qat}: variante \textit{edge} de cerca de 4 bilhões de parâmetros efetivos, com quantização consciente de treinamento. É a menor \textit{build} publicada do E4B ($\sim$5,7~GB) e não cabe em 4~GB de VRAM --- a variante Q4\_K\_M ocupa $\sim$9~GB ---, o que exige \textit{offload} parcial para a CPU na estação de referência;
  \item \textbf{Gemma 4 E2B}~\cite{team2026gemma4}, \textit{tag} \texttt{gemma4:\allowbreak e2b-it-qat}: variante menor da mesma família, com cerca de 2 bilhões de parâmetros efetivos ($\sim$4,0~GB), incluída como alternativa de menor custo computacional ao E4B, que não cabe na VRAM da estação de referência; o E2B também recebe \textit{offload} parcial nessa estação. Permite observar, dentro de uma mesma família, o efeito do tamanho do modelo;
  \item \textbf{Mistral Nemo 12B}~\cite{mistral2024nemo} (Mistral AI e NVIDIA), \textit{tag} \texttt{mistral-nemo:\allowbreak 12b}: 12 bilhões de parâmetros ($\sim$6,6~GB), com suporte consolidado a \textit{Tool Calling} e bom desempenho em português. Representa a faixa superior de tamanho executável com \textit{offload} parcial numa estação modesta.
\end{enumerate}

\subsection{Referência paralela em nuvem}

Como referência --- e não como candidatos à implantação, uma vez que o requisito RNF1 exige execução local ---, a mesma avaliação foi executada sobre \res{groq}{geral}{nmodelos} modelos maiores servidos pelo provedor Groq e disponíveis com suporte a \textit{Tool Calling} em setembro de 2026: \res{groq}{geral}{tags} \cite{qwen2026qwen38, openai2025gptoss}. Os três têm pesos abertos, sob licença Apache 2.0, mas exigem mais memória do que a estação de referência oferece. O objetivo é situar os modelos locais em relação a modelos de maior porte sob o mesmo \textit{prompt}, a mesma ferramenta, o mesmo \textit{dataset} e as mesmas métricas. A troca de provedor usa exatamente a portabilidade descrita no Capítulo~\ref{cap:referencial}: o \texttt{ChatOllama} do LangChain é substituído pelo \texttt{ChatOpenAI} apontado para o \textit{endpoint} compatível do provedor, sem alteração no restante do \textit{pipeline}. Esses resultados são sempre reportados à parte e identificados como execução em nuvem.

% ---
\section{Construção do \textit{Dataset} de Avaliação}
\label{sec:dataset}
% ---

A comparação objetiva dos modelos exige um \textit{dataset} de referência representativo do espectro de consultas que usuários do catálogo tenderiam a formular, cobrindo diferentes níveis de complexidade e combinações de parâmetros, e cujo gabarito seja \textbf{auditável}: cada resposta esperada deve poder ser justificada a partir do texto da consulta e de regras escritas, e qualquer leitor deve poder refazer essa verificação.

O \textit{dataset} serve exclusivamente para \textbf{avaliação}: nenhum modelo foi treinado ou ajustado com ele, e cada modelo recebe as consultas uma a uma, sem exemplos resolvidos (\textit{zero-shot}). A qualidade que importa é, portanto, a do gabarito --- se cada resposta esperada está certa e se as leituras aceitas são todas as leituras legítimas da consulta ---, e ela se apoia em três mecanismos: regras de anotação escritas antes da anotação (Seção~\ref{sec:manual-anotacao}), construção rastreável em camadas (Seção~\ref{sec:camadas}) e uma auditoria com checagens automáticas, anotação independente às cegas, adjudicação registrada e conferência humana de uma amostra (Seção~\ref{sec:auditoria}).

\subsection{Estrutura de cada entrada}

Cada consulta do \textit{dataset} é registrada com os campos da Tabela~\ref{tab:estrutura-dataset}, acrescidos de \texttt{origem} (P, N ou G), \texttt{familia} (P, N, GS \ldots GF) e \texttt{trocas} (leituras estruturais alternativas, das quais \texttt{alternativas} é a expansão); valores alternativos e campos opcionais ficam marcados no próprio \texttt{esperado}. O arquivo \texttt{dataset.json} e o Apêndice~\ref{apendice:dataset} são gerados a partir da mesma fonte, de modo que não podem divergir.

\begin{table}[htbp!]
\centering
\caption{Estrutura de cada entrada do \textit{dataset} de avaliação}
\label{tab:estrutura-dataset}
\small
\begin{tabular}{|l|p{10.5cm}|}
\hline
\textbf{Campo} & \textbf{Descrição} \\
\hline
\texttt{id} & Identificador (P01--P22, N01--N40, GS001 \ldots GF007) \\ \hline
\texttt{fonte} & Rastreabilidade: arquivo e linha de origem (camada P), autores (camada N) ou função geradora, modelo de frase e semente (camada G) \\ \hline
\texttt{categorias} & Uma ou mais das categorias S, C, M, T, O, A; ou F (fora do domínio), que não se combina com as demais \\ \hline
\texttt{consulta} & Texto em linguagem natural \\ \hline
\texttt{esperado} & Leitura preferencial do gabarito: parâmetros de busca na forma canônica, com períodos relativos expressos por regras nomeadas \\ \hline
\texttt{alternativas} & Demais leituras aceitas, quando a consulta admite mais de uma (Seção~\ref{sec:manual-anotacao}) \\ \hline
\texttt{espera\_tool\_call} & Falso quando o comportamento correto é não chamar a ferramenta \\ \hline
\texttt{observacional} & Verdadeiro quando o caso fica fora das métricas principais \\ \hline
\texttt{notas} & Justificativa das leituras alternativas e observações do anotador \\ \hline
\end{tabular}%
\fonte{Elaborada pelos autores.}
\end{table}

\subsection{Categorias de Consultas}

O \textit{dataset} é organizado em seis categorias, alinhadas aos exemplos do protótipo e às demandas do cliente: \textbf{simples} (S, um único parâmetro: ``cartas de São Paulo''); \textbf{compostas} (C, dois ou mais parâmetros: ``ortoimagens de Manaus em 1:50.000''); \textbf{com código MI/INOM} (M: ``INOM SG-22-X-A''); \textbf{com referência temporal relativa} (T: ``nos últimos 3 meses'', ``depois de 2022''); \textbf{com ordenação} (O: ``as 3 cartas mais antigas''); e \textbf{ambíguas ou informais} (A: siglas, grafias sem acento, abreviações e expressões qualitativas). Uma consulta pode pertencer a mais de uma categoria. Uma sétima categoria, \textbf{fora do domínio} (F: ``cartas de Marte'', ``qual a previsão do tempo para amanhã em Manaus?''), reúne as consultas em que a resposta correta é não chamar a ferramenta; ela mede a capacidade de o modelo recusar pedidos que o catálogo não atende.

\subsection{Camadas de Construção}
\label{sec:camadas}

O \textit{dataset} tem 310 consultas, construídas em três camadas complementares (Tabela~\ref{tab:dataset-composicao}):

\begin{enumerate}
  \item \textbf{Camada P (protótipo, 22 casos)}: as consultas do arquivo de testes do protótipo (\texttt{backend/\allowbreak evaluation/\allowbreak test-cases.ts}), cada uma rastreada até a linha de origem. A anotação original da equipe do 1º CGEO é mantida como leitura preferencial --- o \textit{baseline} não é reanotado ---, e as leituras alternativas acrescentadas seguem o manual de anotação, com justificativa registrada. Os gabaritos de datas relativas do arquivo original, calculados para uma data de referência do início de 2025, são substituídos pelas regras nomeadas da Seção~\ref{sec:manual-anotacao}, que se resolvem na data de referência de cada rodada.
  \item \textbf{Camada N (autores, 40 casos)}: redigidos para preencher lacunas de distribuição --- em especial a categoria Simples, ausente na camada P --- e para exercitar ambiguidades e casos de fronteira do domínio.
  \item \textbf{Camada G (geração programática, 248 casos)}: consultas sintetizadas pela combinação de catálogos fechados do domínio --- 27 UFs, 12 municípios de referência, 8 escalas do SCN, 9 tipos de produto das ET-PCDG, 5 CGEOs, 9 projetos, códigos MI e INOM, 15 expressões de tempo (14 relativas e o intervalo absoluto ``entre 2022 e 2023'') e padrões de ordenação --- em modelos de frase do português corrente (com preposição adequada a cada UF e apelidos usuais de escala, como ``25k'' ou ``escala de 25 mil''). O gabarito é conhecido \textbf{por construção}: a frase é montada a partir dos parâmetros que deve produzir. A geração é determinística (semente 42) e rejeita, dentro da camada G, qualquer consulta que repita outra após normalização; a construção do \textit{dataset} verifica ainda que nenhuma consulta gerada repita, após normalização, uma das camadas P e N, e é interrompida se isso ocorrer. A camada inclui sete consultas de fronteira do domínio (família G-F): seis fora do domínio e uma observacional.
\end{enumerate}

\begin{table}[htbp!]
\centering
\caption{Composição do \textit{dataset} de avaliação por origem e família}
\label{tab:dataset-composicao}
\small
\begin{tabular}{|l|>{\centering\arraybackslash}p{1.9cm}|>{\centering\arraybackslash}p{2.3cm}|>{\centering\arraybackslash}p{2.9cm}|}
\hline
\textbf{Origem / família} & \textbf{Consultas} & \textbf{Observa\-cionais} & \textbf{Com mais de uma leitura aceita} \\
\hline
Protótipo do 1º CGEO (P) & 22 & 0 & 11 \\
Autores (N) & 40 & 3 & 15 \\
Geradas --- simples (G-S) & 61 & 0 & 0 \\
Geradas --- compostas (G-C) & 55 & 0 & 3 \\
Geradas --- código MI/INOM (G-M) & 28 & 0 & 0 \\
Geradas --- tempo relativo (G-T) & 41 & 0 & 20 \\
Geradas --- ordenação (G-O) & 19 & 0 & 16 \\
Geradas --- informais (G-A) & 24 & 0 & 2 \\
Geradas --- amplas (G-P) & 13 & 0 & 0 \\
Geradas --- fronteira (G-F) & 7 & 1 & 0 \\
\hline
\textbf{Total} & \textbf{310} & \textbf{4} & \textbf{67} \\
\hline
\end{tabular}
\fonte{Elaborada pelo autor a partir de \texttt{dataset.json}.}
\end{table}


\paragraph{Justificativa da geração programática.} Consultas escritas à mão são caras de redigir e anotar e tendem a variar sobre um pequeno conjunto de padrões familiares aos autores. A geração por modelos de frase sobre catálogos fechados oferece cobertura sistemática de combinações raras (pares UF $\times$ escala pouco frequentes, por exemplo), escala sem custo adicional de anotação e volume suficiente para estimar intervalos de confiança e comparar modelos com testes de significância. A contrapartida é a regularidade linguística, que pode favorecer os modelos: por isso a acurácia é também desagregada por camada (Tabela~\ref{tab:resultados_por_origem}), o que permite distinguir o desempenho nas consultas geradas do desempenho nas consultas redigidas por pessoas (P e N).

\begin{table}[htbp!]
\centering
\caption{Consultas por categoria e por camada (métricas principais; uma consulta pode ter mais de uma categoria)}
\label{tab:dataset-categorias}
\small
\begin{tabular}{|l|c|c|c|c|}
\hline
\textbf{Categoria} & \textbf{P} & \textbf{N} & \textbf{G} & \textbf{Total} \\
\hline
Simples (S) & 0 & 11 & 89 & 100 \\
Compostas (C) & 22 & 21 & 137 & 180 \\
Código MI/INOM (M) & 10 & 5 & 42 & 57 \\
Tempo relativo (T) & 6 & 9 & 45 & 60 \\
Ordenação (O) & 4 & 6 & 19 & 29 \\
Ambíguas/informais (A) & 16 & 19 & 162 & 197 \\
Fora do domínio (F) & 0 & 2 & 6 & 8 \\
\hline
\end{tabular}
\fonte{Elaborada pelo autor.}
\end{table}


\subsection{Manual de Anotação e Leituras Múltiplas}
\label{sec:manual-anotacao}

O gabarito é construído segundo um manual de anotação escrito (Apêndice~\ref{apx:manual}) que fixa seis princípios --- evidência explícita no texto, forma canônica, leituras múltiplas explícitas, tratamento das consultas fora do domínio, consistência e casos observacionais --- e regras para cada um dos doze campos da ferramenta.

O princípio mais importante para a confiabilidade da avaliação é o das \textbf{leituras múltiplas explícitas}. Uma parte das consultas admite, legitimamente, mais de uma resposta correta: ``cartas de São Paulo'' pode referir-se ao estado ou à capital; ``produtos publicados nos últimos 3 meses'' pode significar noventa dias ou três meses de calendário; ``as cartas mais recentes'' pode ser ordenado pela data de publicação ou pela de criação, pois a consulta não diz qual. Um gabarito que escolhesse uma única leitura penalizaria como erro uma resposta razoável e mediria a coincidência com a preferência do anotador, não a qualidade da tradução. Nesses casos o gabarito lista \textbf{todas} as leituras aceitas, por três mecanismos: \emph{valor alternativo} (o campo aceita mais de um valor), \emph{campo opcional} (pode faltar sem erro, como \texttt{limit = 1} em ``a carta mais recente'') e \emph{leitura estrutural alternativa} (o mesmo trecho pode ocupar campos diferentes, como \texttt{state} ou \texttt{city}). As expressões de tempo relativo são registradas como regras nomeadas, resolvidas na data de referência da rodada --- registrada no manifesto e informada ao modelo no \textit{prompt} de sistema --- com as leituras aceitas do Quadro~\ref{quadro:tempo-relativo}.

A ambiguidade não é, porém, tratada com leniência genérica: toda leitura aceita está escrita no gabarito, justificada por uma regra do manual e visível no Apêndice~\ref{apendice:dataset}; fora delas, a comparação é estrita (Seção~\ref{sec:metricas}). Das 306 consultas das métricas principais, a Tabela~\ref{tab:dataset-composicao} indica quantas admitem mais de uma leitura.

\paragraph{Consultas fora do domínio e casos observacionais.} Oito consultas estão fora do domínio (N36, N37 e seis da família G-F): o comportamento correto é determinado --- não chamar a ferramenta --- e elas entram nas métricas principais, na categoria F. Quatro consultas são observacionais (N38, N39, N40 e GF007), porque o comportamento correto não é determinável a partir do texto: valor impossível de representar (``escala 1:10.000.000'', ``cartas do estado 42'') ou pedido de resposta discutível (``cartas topo no oceano atlântico'', ``cartas do futuro (ano 2050)''). Elas são executadas e o comportamento de cada modelo é reportado à parte; as métricas principais cobrem as 306 consultas restantes.

\begin{table}[htbp!]
\centering
\caption{Frequência de cada campo na leitura preferencial do gabarito (métricas principais)}
\label{tab:dataset-campos}
\small
\begin{tabular}{|l|c||l|c|}
\hline
\textbf{Campo} & \textbf{Ocorr.} & \textbf{Campo} & \textbf{Ocorr.} \\
\hline
\texttt{keyword} & 62 & \texttt{project} & 11 \\
\texttt{scale} & 84 & \texttt{publicationPeriod} & 51 \\
\texttt{productType} & 44 & \texttt{creationPeriod} & 10 \\
\texttt{state} & 130 & \texttt{sortField} & 29 \\
\texttt{city} & 27 & \texttt{sortDirection} & 29 \\
\texttt{supplyArea} & 40 & \texttt{limit} & 21 \\
\hline
\end{tabular}
\fonte{Elaborada pelos autores.}
\end{table}


% ---
\section{Auditoria do \textit{Dataset}}
\label{sec:auditoria}
% ---

A confiabilidade das métricas depende diretamente da confiabilidade do gabarito. Por isso o \textit{dataset} foi submetido a uma auditoria em três etapas, repetidas até não restar divergência sem decisão registrada. O procedimento, as ferramentas e os resultados completos estão no Apêndice~\ref{apx:auditoria} e no repositório (\texttt{data/auditoria/}).

\begin{enumerate}
  \item \textbf{Checagens automáticas} (\texttt{pfc-auditar}): unicidade de identificadores e de consultas (após normalização de caixa, acentos e pontuação); validade de todos os valores contra os enumerados da ferramenta; resolução de todas as regras de tempo relativo em datas diversas, inclusive anos bissextos e viradas de ano; rastreabilidade das consultas da camada P até a linha citada do arquivo de origem; um \textbf{anotador por regras}, que localiza na consulta o trecho que justifica cada campo do gabarito e aponta campo sem evidência, evidência sem campo e valor divergente; e uma verificação de \textbf{consistência} de convenções entre casos (a mesma expressão, a mesma anotação).
  \item \textbf{Anotação independente às cegas}: todas as 310 consultas foram anotadas novamente, a partir do zero, por um modelo de linguagem de propósito geral de grande porte (Claude Opus 5.5, da Anthropic), que não integra o conjunto de modelos avaliados. A anotação foi feita em dez lotes de 31 consultas, em sessões independentes, cada uma com acesso apenas ao manual de anotação, à definição da ferramenta, à data de referência e ao próprio lote --- sem acesso ao gabarito, ao gerador ou aos demais lotes. Para cada consulta, o anotador registrou a decisão (chamar a ferramenta, não chamar ou caso observacional), a leitura preferencial, as leituras alternativas previstas pelo manual e uma justificativa por referência às regras.
  \item \textbf{Adjudicação}: cada consulta com divergência em qualquer das medidas de concordância recebeu uma decisão registrada --- gabarito mantido, corrigido ou ampliado com leitura alternativa --- com justificativa por referência ao manual. O gabarito anterior a cada decisão fica registrado, de modo que a concordância anterior à adjudicação pode ser recalculada a partir do repositório.
\end{enumerate}

A concordância entre o gabarito e a anotação independente foi medida em rigor crescente (Tabela~\ref{tab:auditoria}): a decisão sobre chamar a ferramenta, com o kappa de Cohen~\cite{cohen1960coefficient,artstein2008intercoder}, que desconta a concordância esperada ao acaso; a compatibilidade das leituras nos dois sentidos (a leitura do anotador é aceita pelo gabarito e a leitura preferencial do gabarito é aceita pelo anotador); a identidade da leitura preferencial; a presença e o valor de cada campo; e a detecção de ambiguidade, isto é, se a consulta admite mais de uma leitura. A coluna ``antes da adjudicação'' mede a concordância propriamente dita; a coluna ``após'' mostra o que restou depois das decisões.

\begin{table}[htbp!]
\centering
\caption{Resultados da auditoria do \textit{dataset}}
\label{tab:auditoria}
\small
\begin{tabular}{|p{7.3cm}|c|c|}
\hline
\textbf{Verificação} & \textbf{Antes da adjudicação} & \textbf{Após} \\
\hline
\multicolumn{3}{|l|}{\textit{Checagens automáticas}} \\ \hline
Problemas estruturais (ids, duplicatas, enumerados, datas) & --- & 0 \\ \hline
Consultas da camada P conferidas na linha de origem & --- & 22/22 \\ \hline
Apontamentos do anotador por regras sem resolução & --- & 0 \\ \hline
Expressões com anotação inconsistente sem justificativa & --- & 0 \\ \hline
\multicolumn{3}{|l|}{\textit{Anotação independente às cegas (310 consultas)}} \\ \hline
Decisão (chamar / não chamar / observacional) & 302/310 (97,4\%) & 310/310 (100,0\%) \\
\quad kappa de Cohen & 0,658 & 1,000 \\ \hline
Leituras compatíveis nos dois sentidos & 300/310 (96,8\%) & 309/310 (99,7\%) \\ \hline
Leitura preferencial idêntica$^{a}$ & 294/298 (98,7\%) & 294/298 (98,7\%) \\ \hline
Presença de cada campo na leitura preferencial$^{a,b}$ (kappa) & 0,999 & 0,999 \\ \hline
Valor do campo, quando ambos o anotam$^{a}$ & 534/537 (99,4\%) & 534/537 (99,4\%) \\ \hline
Consulta com mais de uma leitura aceita$^{a}$ & 296/298 (99,3\%) & 297/298 (99,7\%) \\
\quad kappa de Cohen & 0,981 & 0,990 \\ \hline
Consultas com divergência em alguma medida, sem decisão registrada & 14 & 0 \\ \hline
\end{tabular}
\fonte{Elaborada pelo autor a partir de \texttt{data/auditoria/}. (a) Sobre as consultas em que gabarito e anotador pedem a chamada da ferramenta nas métricas principais. (b) Doze campos por consulta.}
\end{table}


As medidas de campo e de ambiguidade mostram concordância quase perfeita já antes da adjudicação. A principal divergência foi de classificação: o gabarito tratava como observacionais oito consultas fora do domínio cujo comportamento correto é determinado pelo próprio manual (não chamar a ferramenta). A decisão foi corrigir o gabarito e incluí-las nas métricas principais, na categoria F. As demais divergências foram uma leitura ampliada (N33, em que ``cartografia sistemática'' passou a admitir a leitura sem o campo \texttt{project}, pois a expressão não consta da descrição da ferramenta) e cinco divergências mantidas com justificativa, entre elas três escolhas de leitura preferencial dentro do mesmo conjunto de leituras aceitas, que não afetam a pontuação. O Quadro~\ref{quadro:adjudicacao} do Apêndice~\ref{apx:auditoria} lista essas 14 decisões e mais uma, originada nas checagens de consistência (P14, mantida), num total de 15 consultas adjudicadas.

Duas ressalvas delimitam o alcance da auditoria. O anotador por regras foi escrito pelos autores e, portanto, verifica evidência e consistência, mas não constitui uma anotação independente. A anotação às cegas é independente do gabarito, mas foi feita por um modelo de linguagem, e não por uma pessoa: ela identifica divergências com alta sensibilidade, e a decisão sobre cada uma é tomada na adjudicação, por referência ao manual. Na camada G, além disso, o gabarito é conhecido por construção, e a auditoria verifica sobretudo se a frase gerada sustenta, em português corrente, os parâmetros que a geraram.

Para compensar em parte essa ressalva, uma amostra estratificada de 40 consultas (semente 42), sorteada entre as camadas P, N e G, foi conferida manualmente por um dos autores, que confirmou o gabarito de todas elas (Tabela~\ref{tab:auditoria}). Em síntese, a resposta à pergunta ``como saber se o gabarito está correto'' é dada por números: antes de qualquer correção, uma anotação feita do zero, sem acesso ao gabarito, concordou com ele em 99,4\% dos valores de campo anotados por ambos (534 de 537), teve kappa de 0,999 na presença dos campos e de 0,981 na identificação das consultas ambíguas, e encontrou leitura compatível em 300 das 310 consultas; a única divergência sistemática --- a classificação de oito consultas fora do domínio --- foi corrigida na adjudicação, e nenhuma divergência ficou sem decisão registrada.

% ---
\section{Lote de Validação Independente}
\label{sec:lote-validacao}
% ---

As 310 consultas deixam três perguntas que só um segundo conjunto responde: se as conclusões dependem da regularidade da camada G; se a recusa, medida sobre oito consultas fora do domínio, se sustenta numa amostra maior; e se o desempenho do \textit{Tool Calling} reflete a abordagem ou a especificação dada ao modelo. Para isso foi construído um \textbf{lote de validação} (daqui em diante, também \textbf{lote 1}) de \res{lotedados}{geral}{aceitas} consultas novas, independente das 310, usado no Capítulo~\ref{cap:resultados} com o modelo local de maior acurácia com \textit{Tool Calling}. O lote difere das 310 em três pontos de método: os valores vêm de catálogos reais; o gabarito de cada consulta é fixado \emph{antes} de o texto existir; e a consulta só entra no lote se uma anotação às cegas, feita sem acesso ao gabarito, for compatível com ele. Cada etapa tem \textit{script}, semente e arquivos versionados no repositório, e a metodologia completa está em \texttt{docs/lote\_validacao.md}.

\paragraph{Fontes reais.} Os metadados públicos do catálogo do BDGEx foram coletados pelo serviço CSW (OGC CSW 2.0.2) em 28/09/2026: 31.190 registros, dos quais saem nomes de folha (7.838 distintos), códigos MI e INOM, escalas, tipos de produto, CGEOs e anos. A lista oficial de municípios do IBGE (5.571) fornece os municípios. Nas consultas compostas, metade das combinações de escala, tipo, CGEO e folha é tirada de uma mesma carta do acervo; o lugar é sorteado à parte.

\paragraph{Alvos e gabarito por construção.} Foram sorteados \res{lotedados}{geral}{alvos} \emph{alvos} (semente 2026) em oito famílias (Tabela~\ref{tab:lote_composicao}). Cada alvo especifica o gabarito, segundo o mesmo manual de anotação, a forma de superfície de cada valor (``25k'', ``1:25000'', ``3o cgeo'', ``terceiro cgeo'', sigla ou nome da UF, INOM compacto em minúsculas\ldots), o que a consulta não pode dizer, um registro de linguagem (formal, direto, coloquial, pergunta, telegráfico, sem acento, com erro de digitação ou com contexto de uso) e uma persona (do oficial que planeja uma operação ao cidadão sem conhecimento técnico). O manual recebeu extensões que não alteram nenhuma anotação das 310 (Apêndice~\ref{apx:manual}, ``Extensões para o lote de validação''): regras de tempo relativo parametrizadas (``últimos $N$ meses'', ``desde AAAA'' para qualquer $N$ e ano), municípios acompanhados da sigla da UF e as \emph{armadilhas de vocabulário} (``carta de vinhos'', ``folha de pagamento'', ``escala de plantão''), que estão fora do domínio.

Duas mudanças de composição respondem diretamente às perguntas acima. A categoria F passa a ter \res{lotedados}{geral}{nF} consultas, em nove subtipos, o que permite medir a recusa como classificação --- precisão, \textit{recall} e F1 --- e a taxa de falsa recusa nas consultas do domínio. E uma nova categoria, \textbf{subespecificadas} (E, \res{lotedados}{geral}{nE} consultas), reúne pedidos de produtos do acervo sem nenhum critério representável no \textit{schema} (``quero ver as cartas do acervo'', ``cartas do Nordeste'', ``mapas para uma trilha''), em que duas respostas são aceitas: buscar sem filtros ou não buscar. Como prevê o manual (``pedir esclarecimento ou recusar''), a pontuação aceita qualquer não busca, e não só o pedido de esclarecimento: a recusa de uma consulta E também conta como acerto, embora a especificação v2, descrita adiante, diga que ``consultas vagas sobre o acervo não são fora do escopo''. A regra favorece sobretudo a Saída Estruturada v2, que marcou \texttt{fora\_do\_escopo} em \res{lote}{se2}{enao} das consultas E; com essas recusas contadas como erro, a sua acurácia no lote seria de \res{lote}{se2}{accEerro}, e não de \res{lote}{se2}{acuracia} (Seção~\ref{sec:lote-resultados}). Um único critério (``ortoimagens'') continua sendo consulta simples, que exige a busca, como nas 310. Ao todo, \res{lotedados}{geral}{multiplas} consultas do lote admitem mais de uma leitura aceita.

\paragraph{Redação controlada.} Os alvos foram embaralhados e divididos em 15 tarefas de 130, cada uma redigida por um agente de linguagem (Claude Opus 5.5, da Anthropic, que não integra os modelos avaliados), que recebeu apenas o que a consulta deveria e não deveria dizer, o registro e a persona --- não o gabarito --- e a instrução de soar como uma pessoa digitando numa caixa de busca, variando a construção de uma mensagem para outra. Checagens automáticas conferem, em cada consulta, a presença literal das superfícies exigidas, o verbo do período (publicação, criação ou nenhum) e a ausência de marcadores de campos que o alvo não tem; a deduplicação descarta consultas que repitam, após normalização, uma das 310 ou outra do lote.

\paragraph{Anotação às cegas e filtro.} Todas as consultas redigidas foram anotadas de novo, a partir do zero, por outros 15 agentes, com o mesmo formato e as mesmas referências da anotação independente das 310 (Seção~\ref{sec:auditoria}): o manual, a definição da ferramenta, a data de referência e a lista de municípios do IBGE. As consultas chegaram embaralhadas e com identificadores neutros, sem família, persona ou gabarito. Uma amostra aleatória de 20\% (\res{lotedados}{geral}{nB} consultas) recebeu uma segunda anotação, por outros três agentes, que não participa do filtro. Entra no lote a consulta cuja anotação é compatível com o gabarito nos dois sentidos --- o critério da auditoria das 310. Das \res{lotedados}{geral}{redigidas} consultas redigidas, \res{lotedados}{geral}{aceitas} (\res{lotedados}{geral}{pctaceitas}) formam o lote; as descartadas, com o motivo, estão na Tabela~\ref{tab:lote_composicao} e no repositório.

\begin{table}[htbp!]
\centering
\caption{Lote de validação: do alvo ao lote final, por família}
\label{tab:lote_composicao}
\footnotesize
\begin{tabular}{|l|l|c|c|c|c|}
\hline
\textbf{Família} & \textbf{Testa} & \textbf{Alvos} & \textbf{Checagens} & \textbf{Anotação às cegas} & \textbf{Lote final} \\
\hline
VS & Simples (um critério) & 280 & $-$3 & 0 & 277 \\
VC & Compostas & 360 & 0 & $-$1 & 359 \\
VM & Códigos MI/INOM & 160 & 0 & 0 & 160 \\
VT & Tempo & 260 & $-$1 & 0 & 259 \\
VO & Ordenação & 200 & $-$1 & 0 & 199 \\
VA & Leituras múltiplas & 200 & $-$7 & 0 & 193 \\
VE & Subespecificadas & 150 & $-$4 & 0 & 146 \\
VF & Fora do domínio & 340 & $-$4 & 0 & 336 \\
\hline
\textbf{Total} & & 1.950 & $-$20 & $-$1 & \textbf{1.929} \\
\hline
\end{tabular}
\fonte{Elaborada pelos autores. Checagens: descartadas nas checagens automáticas e na deduplicação. Anotação às cegas: descartadas por incompatibilidade entre o gabarito por construção e a anotação A.}
\end{table}


A concordância, medida sobre todas as consultas redigidas antes do filtro (Tabela~\ref{tab:lote_concordancia}), foi de kappa \res{lotedados}{geral}{kappadecisao} na decisão (buscar, buscar ou esclarecer, não buscar), \res{lotedados}{geral}{compat} de leituras compatíveis e \res{lotedados}{geral}{valoresiguais} dos valores de campo; entre os dois anotadores, \res{lotedados}{geral}{compatAB} de leituras compatíveis. Esses números são altos por construção --- os alvos exigem superfícies explícitas, e o manual decide as ambiguidades previstas --- e devem ser lidos com a mesma ressalva da auditoria das 310, agravada: redatores e anotadores são modelos da mesma família. Eles mostram que cada consulta sustenta o seu gabarito segundo o manual, de forma reproduzível por anotadores independentes; não mostram que pessoas leriam do mesmo modo. Dois sinais indicam que a anotação não reproduziu o gabarito por outra via: na amostra dupla, os anotadores concordam mais entre si do que com a construção, e as justificativas que escreveram são textualmente distintas. O texto das consultas é sintético --- não há registro de consultas reais de usuários do BDGEx disponível ---, e essa é a principal limitação do lote.

\begin{table}[htbp!]
\centering
\caption{Concordância na anotação às cegas do lote de validação}
\label{tab:lote_concordancia}
\footnotesize
\begin{tabular}{|l|c|c|}
\hline
\textbf{Medida} & \textbf{Construção $\times$ A} (n=1.950) & \textbf{A $\times$ B} (n=390) \\
\hline
Decisão (buscar / buscar ou esclarecer / não buscar), kappa & 1,000 & 1,000 \\
Leituras compatíveis nos dois sentidos & 99,8\% & 100,0\% \\
Leitura preferencial idêntica & 99,8\% & 100,0\% \\
Presença dos campos, kappa & 0,999 & 1,000 \\
Valor igual quando ambos preenchem & 100,0\% & 100,0\% \\
Detecção de leituras múltiplas, kappa & 0,967 & 1,000 \\
\hline
\end{tabular}
\fonte{Elaborada pelos autores. Medidas de \texttt{audit.concordancia}, as mesmas da auditoria das 310 consultas, sobre todas as consultas redigidas (antes do filtro). A: anotação de todas as consultas; B: segunda anotação independente de uma amostra aleatória de 20\%.}
\end{table}


\paragraph{Especificação v2.} A análise de erros do \textit{Tool Calling} do Gemma 4 E4B nas 310 consultas mostrou que os dois erros mais frequentes contrariam instruções que o \textit{prompt} v1 já dava. Em \res{diag}{tc1}{naochamou} das \res{diag}{tc1}{execucoes} execuções (\res{diag}{tc1}{pctnaochamou}), o modelo não buscou numa consulta do domínio, embora a instrução 1 diga que ``um único parâmetro basta'' e a 2 só permita deixar de chamar a ferramenta quando a consulta não tratar de produtos do acervo; em \res{diag}{tc1}{naochamouesclarece} delas pediu esclarecimento --- às vezes de informação que já tinha, como ``preciso saber qual é o ano atual''. Em \res{diag}{tc1}{tipoinventado}, acrescentou o tipo de produto a ``cartas'' ou ``mapas'', contra a instrução 3 (``não complete campos por suposição''). Outros vêm de exemplos que a descrição trazia ou de convenções do manual que ela não informava: em \res{diag}{tc1}{codigoalterado}, o modelo acrescentou ao código um sufixo, em geral o do exemplo da descrição (``2901'' virou ``2901-2-NE''), ou usou ``folha'' como \textit{keyword}; em \res{diag}{tc1}{periodohoje}, resolveu ``este ano'' como um período que começa hoje. Em \res{diag}{tc1}{escalasemponto}, deixou a escala sem o ponto de milhar, embora o enumerado da definição traga a grafia com o ponto. A descrição v1 também não informava outras convenções do gabarito: que ``última atualização'' indica a data de criação, que um período que vai até hoje tem fim e que uma escala qualitativa como ``grande escala'' deve ser preenchida com uma das escalas que ela abrange.

Para separar o efeito da abordagem do efeito da especificação, as duas abordagens foram avaliadas também numa \textbf{v2}, que muda apenas a especificação: descrições dos parâmetros que explicitam as convenções do manual (quando preencher cada campo, a forma canônica, o que não deduzir), sem os códigos de exemplo da v1 e com a lista das siglas das 27 UFs; no \textit{Tool Calling}, uma segunda ferramenta, \texttt{recusar\_consulta}, e a instrução de chamar exatamente uma das duas, sem pedir esclarecimento; na Saída Estruturada, as mesmas descrições, a mesma proibição de pedir esclarecimento e um campo \texttt{fora\_do\_escopo}, que torna a recusa representável. A v2 não ficou livre de exemplos copiáveis: a descrição de \texttt{keyword} traz ``sf22yd'' → ``SF-22-Y-D'', que é o código da consulta N35 das 310 com o seu gabarito, e a de \texttt{limit}, ``a carta mais recente'' → 1. A descrição de \texttt{state} pede o nome por extenso, mas troca a regra com direção da v1 (``Normalizar siglas: `RJ' → `Rio de Janeiro'\,'') por uma lista de pares (``AC Acre, AL Alagoas'', e assim por diante), e os dois \textit{prompts} v2 perderam a frase da instrução 1 da v1 segundo a qual siglas, abreviações e grafias sem acento ``são esperadas e devem ser reconhecidas e normalizadas''. Tipos, enumerados e mecanismos não mudam, e as duas v2 recebem as mesmas descrições dos 12 parâmetros; o texto que enquadra a recusa, porém, difere. Só o \textit{Tool Calling} é instruído a buscar ``mesmo que de forma vaga'', lê na descrição da busca que ``uma consulta sem nenhum parâmetro reconhecível ainda é uma busca'' e, na da ferramenta de recusa, que ela não serve para consultas vagas sobre o acervo; na Saída Estruturada, recusar é responder \texttt{fora\_do\_escopo} igual a \texttt{true} ``e nada mais'', e a descrição do campo diz que ``consultas vagas sobre o acervo não são fora do escopo''. A comparação entre as duas v2 mede, portanto, o mecanismo junto com esse texto, sobretudo na recusa.

Como a v2 foi desenhada a partir dos erros nas 310, esperava-se que o seu resultado nesse conjunto, de dentro da amostra, fosse otimista. Não foi: nas 310, o \textit{Tool Calling} v2 acertou \res{base}{tc2}{acuracia} e a Saída Estruturada v2, \res{base}{se2}{acuracia}, abaixo das v1 no mesmo conjunto (\res{base}{tc1}{acuracia} e \res{base}{se1}{acuracia}). No domínio, a diferença entre a v2 e a v1 foi, no \textit{Tool Calling}, de \res{base}{dominio}{tcv2menosv1} p.p. nas 310 e de \res{lote}{dominio}{tcv2menosv1} p.p. no lote e, na Saída Estruturada, de \res{base}{dominio}{sev2menosv1} e \res{lote}{dominio}{sev2menosv1} p.p.: dentro da amostra, a v2 foi pior do que fora dela no \textit{Tool Calling} e parecida na Saída Estruturada. Das consultas do domínio que o \textit{Tool Calling} v1 acertava, o v2 errou \res{base}{tc1tc2}{taxaquebra} nas 310, contra \res{lote}{tc1tc2}{taxaquebra} no lote, em boa parte porque os erros de forma que a v2 introduziu, como a sigla no lugar do nome da UF e o prefixo ``MI'' na \textit{keyword}, atingem formas que se repetem na camada G, como ``cartas de'' seguido de uma UF e ``folha MI'' seguida do código e da UF. Os consertos dentro da amostra existem, mas são menores que as quebras: nas 310, o \textit{Tool Calling} v2 acertou \res{base}{tc1tc2}{frconsertadas} consultas em que o v1 não buscava e \res{base}{tc1tc2}{camposconsertados} em que o v1 buscava com erro, mas errou \res{base}{tc1tc2}{quebradas} que o v1 acertava. A medida que vale continua sendo a do lote, que o desenho da v2 não viu (Seção~\ref{sec:lote-resultados}).

\paragraph{Execução e medidas.} O Gemma 4 E4B foi executado na mesma T4 das rodadas completas, com uma repetição, nas quatro configurações (\textit{Tool Calling} e Saída Estruturada, v1 e v2) sobre o lote e nas duas v2 sobre as 310. Além das métricas da Seção~\ref{sec:metricas}, o lote permite medir a recusa como classificação (positivo: consulta fora do domínio; previsto positivo: não buscou), o comportamento na categoria E, a acurácia nas consultas com uma e com mais de uma leitura aceita e a acurácia por registro de linguagem. As comparações entre configurações usam o teste de McNemar exato e, para a diferença de acurácia e de F1, um \textit{bootstrap} pareado com 2.000 reamostragens das consultas.

% ---
\section{Especificação v3: Ferramentas Auxiliares e Validação com Retorno}
\label{sec:v3}
% ---

Os resultados do lote de validação (Seção~\ref{sec:lote-resultados}) mostraram que os erros do \textit{Tool Calling} com a v2 são, em boa parte, de forma --- sigla no lugar do nome da UF, prefixo no código, limite em plurais --- e reversíveis por código, e que o \textit{Tool Calling} recusa com precisão muito maior que a Saída Estruturada. A v3 pergunta o que o \textit{Tool Calling} ganha quando se usa o que ele tem de próprio: chamar funções no meio da resposta e receber um retorno delas. Ela acrescenta à v2, em degraus, as peças abaixo, descritas em \texttt{docs/v3.md} no repositório.

\paragraph{Descrições corrigidas.} As descrições dos parâmetros são as da v2 com os efeitos colaterais medidos no lote 1 corrigidos: a sigla aparece só como entrada (``MT'' $\rightarrow$ ``Mato Grosso''), o limite no singular exige ordenação, o código vai sem prefixo e sem sufixo acrescentado, o município vai com o nome inteiro, e o tipo de produto só é preenchido com evidência na consulta (``na dúvida, não preencha'').

\paragraph{Ferramentas auxiliares.} Quatro funções determinísticas que o modelo pode chamar antes da busca: \texttt{identificar\_nome}, que diz se um trecho da consulta é UF (inclusive sigla), município (pela lista do IBGE, com a UF), folha do acervo (pelo índice do BDGEx) ou região, e dá a forma canônica; \texttt{normalizar\_codigo}, que devolve o código MI ou INOM sem prefixo e com hífens e diz se ele consta no acervo; \texttt{normalizar\_escala}, que converte a menção à escala para o enumerado do SCN; e \texttt{resolver\_periodo}, que resolve uma expressão de tempo pela tabela do manual, tirando a aritmética de datas do modelo. A elas se soma \texttt{pedir\_esclarecimento}, que o \textit{prompt} e a descrição da ferramenta reservam às consultas sem nenhum critério (o código não impõe a restrição), com um usuário simulado de resposta fixa: não tem outros detalhes, e a busca pode seguir com o que foi dito. Uma pergunta que o modelo escreve como texto recebe a mesma resposta, como receberia numa conversa, e uma resposta só em texto, sem pergunta nem chamada, recebe uma vez o pedido ``Responda chamando uma ferramenta''. Chamadas que o modelo escreve como texto, em vez de emiti-las como \textit{tool call} --- na sintaxe de função, em JSON ou no formato nativo do Gemma ---, são interpretadas como chamadas e contadas em todas as configurações v3, inclusive na de uma chamada; na v1 e na v2, a mesma resposta conta como não busca. Essa leitura é parte do sistema medido e teria de acompanhar a v3 numa implantação. Cada consulta tem no máximo seis chamadas ao modelo, três tentativas de busca e duas perguntas.

\paragraph{Retorno: validação e recusa contestada.} Antes de executar, a ferramenta de busca confere os parâmetros e devolve ao modelo, como resultado da chamada, os erros de forma (sigla no lugar do nome, prefixo no código, escala fora do enumerado, UF ou palavra genérica como \textit{keyword}), que impedem a busca, e avisos de evidência: campo preenchido sem trecho correspondente na consulta (tipo de produto a partir de ``cartas'', limite sem número, estado não citado, período que não corresponde à expressão pela tabela do manual) e campo ausente cuja evidência está na consulta (UF ou município citado, código, expressão de tempo). Depois de um aviso, repetir os mesmos parâmetros confirma a busca: o aviso orienta, não obriga. A recusa também tem retorno: se a consulta cita um critério forte do catálogo --- tipo de produto, código, escala, CGEO, projeto, UF ou município ---, a primeira recusa é devolvida ao modelo com a lista desses critérios e a lembrança de que a finalidade declarada pelo usuário (uma obra, um trabalho) não impede a busca; uma segunda recusa vale. A regra contestaria \res{v3}{recusa}{contestadasF} das \res{v3}{recusa}{nF} consultas fora do domínio do lote 1 (``o que é um MDT'', ``PIB da Bahia''), que o modelo precisa então confirmar.

Como o validador codifica as convenções do manual, ele foi testado contra as respostas corretas: aplicado à leitura preferencial do gabarito de todas as consultas de desenvolvimento, não pode reclamar de nenhuma. A primeira versão reclamava de \res{v3}{autoteste}{inicial} das \res{v3}{autoteste}{n} respostas certas do lote 1; a versão congelada, de \res{v3}{autoteste}{final} no lote 1 e de \res{v3}{autoteste}{final310} das \res{v3}{autoteste}{n310} nas 310. A mesma codificação é a principal ressalva da v3: as ferramentas e o validador implementam as convenções que definem o gabarito --- para datas, a mesma tabela que resolve as regras de tempo ---, de modo que o resultado mede o conjunto do modelo com o código do domínio, e não só o modelo. Por isso a análise separa a primeira decisão --- a primeira busca ou recusa, antes de qualquer retorno do validador --- da resposta final. No \textit{Tool Calling}, essa primeira decisão já vem depois das ferramentas auxiliares, que também são código do domínio: o que elas acrescentam aparece nos degraus, e o que o modelo faz só com a especificação é medido pelo degrau das descrições corrigidas.

\paragraph{Degraus e controle.} Cada configuração acrescenta uma peça à anterior: o \textit{Tool Calling} v1 (a solução do Capítulo~\ref{cap:aprimoramento}); a v2; a v3 só com as descrições corrigidas, ainda numa chamada; com as ferramentas auxiliares, o pedido de esclarecimento com usuário simulado e o pedido para responder chamando uma ferramenta, num laço de chamadas; e com o retorno, que é a v3 completa. Os degraus são cumulativos e medidos numa só ordem: cada diferença é o efeito da peça dadas as anteriores. O controle é a Saída Estruturada com as mesmas descrições v3, o mesmo validador e a mesma recusa contestada, aplicados como nova tentativa, do modo como o protótipo fazia com o Instructor. Ele não recebe as ferramentas auxiliares, o usuário simulado, o pedido para responder chamando uma ferramenta nem a leitura de chamadas escritas como texto, e o seu \textit{prompt} não traz a instrução, dada só ao \textit{Tool Calling}, de que nome de lugar desconhecido não é motivo de recusa. Se ele empatar com a v3, o ganho da v3 é das descrições e do validador, que os dois recebem; se ficar abaixo, a diferença mede o mecanismo de \textit{Tool Calling} junto com as ferramentas auxiliares, que só ele chama no meio da resposta, e com as demais peças que só ele tem, sem separá-los. Um controle que desse à Saída Estruturada, antes da resposta, o que as ferramentas auxiliares devolvem separaria o acesso a esse código do mecanismo; ele não foi executado. As referências externas são a melhor configuração v1 ou v2 isolada no lote 2 e a arquitetura em duas etapas simulada na Seção~\ref{sec:lote-resultados} (o \textit{Tool Calling} v2 decide se busca e a Saída Estruturada v1 extrai os parâmetros), calculada no lote 2 com as rodadas das duas configurações.

\paragraph{Desenvolvimento, teste e pré-registro.} As ferramentas, as descrições e o validador foram desenvolvidos só com as 310 consultas e o lote 1, numa amostra estratificada de 160 consultas do lote 1 executada no Gemma 4 E4B local (o mesmo arquivo de modelo da T4), em ciclos registrados no repositório (\texttt{docs/v3\_desenvolvimento.md}). Como cada ciclo corrigia erros observados nessas mesmas consultas, os números de desenvolvimento são de dentro da amostra e só são citados como registro do processo, não como resultado. Uma variante com orientações escritas a partir dos erros de desenvolvimento --- regras curtas em arquivos \textit{markdown}, injetadas no \textit{prompt} quando pertinentes à consulta, na linha das memórias de lições de agentes \cite{zhao2024expel} --- foi testada e descartada no desenvolvimento, antes do teste, por não ter efeito mensurável: com e sem elas, a v3 acertou o mesmo número de consultas, \res{v3dev}{orientacoes}{sem} das \res{v3dev}{orientacoes}{n} (\res{v3dev}{orientacoes}{socom} acertadas só com elas e \res{v3dev}{orientacoes}{sosem} só sem elas), e a primeira decisão foi um pouco melhor sem elas. Retirá-las simplifica o desenho, mas não o livra do risco de sobreajuste: as regras do validador também foram escritas em ciclos sobre os erros dessas mesmas consultas, e é o lote 2 que mede quanto isso pesa.

O teste é um \textbf{lote 2}, construído com a mesma receita do lote 1 --- outra semente e outros agentes redatores e anotadores, da mesma família de modelos ---, com \res{lote2dados}{geral}{aceitas} consultas (\res{lote2dados}{geral}{nF} fora do domínio e \res{lote2dados}{geral}{nE} subespecificadas; concordância na decisão com kappa de \res{lote2dados}{geral}{kappadecisao} e \res{lote2dados}{geral}{compat} de leituras compatíveis) e nenhuma delas lida no desenvolvimento \cite{dwork2015reusable}; os relatórios dos anotadores do lote 2 citaram trechos de algumas dezenas de consultas, e nada da v3 foi escrito ou alterado a partir deles. Uma checagem automática confirmou que nenhum exemplo das descrições, das ferramentas ou dos \textit{prompts} coincide com uma consulta inteira do lote 2; o que coincide é vocabulário do domínio, como ``100k'' e ``mais recente''. Como os dois lotes sortearam nomes e códigos das mesmas listas que as ferramentas consultam (a lista de municípios do IBGE e o índice do BDGEx) e exigem na consulta a superfície literal de cada valor, o teste não mede o alcance das ferramentas em nomes e códigos fora dessas listas. Antes da rodada de teste, a v3 foi congelada: o caderno do Colab grava o \textit{hash} SHA-256 dos \textit{prompts}, das ferramentas e do código e dos dados das ferramentas, e se recusa a rodar com outra versão; o plano de análise --- a comparação principal com o controle, os degraus, as referências e a decomposição entre primeira decisão e resposta final --- foi fixado e versionado no repositório (\texttt{docs/v3.md}) antes da rodada; qualquer outra análise do lote 2, como a leitura dos erros e a acurácia por categoria, é exploratória, foi feita depois do teste e não alterou a v3.

% ---
\section{Protocolo de Avaliação}
\label{sec:protocolo}
% ---

\subsection{Visão Geral}

Cada modelo executa todas as consultas do \textit{dataset}, em ordem fixa, sob condições controladas e reproduzíveis. A execução é automatizada pelo comando \texttt{pfc-avaliar} (Capítulo~\ref{cap:aprimoramento}), que grava um registro por consulta e repetição --- consulta, gabarito resolvido, parâmetros emitidos pelo modelo, \textit{tool calls} brutos, texto de resposta, latência, contagem de \textit{tokens} e classe de erro --- e um manifesto de reprodutibilidade por rodada.

\subsection{Ambientes de Execução}
\label{sec:ambiente}

\paragraph{Estação de referência.} A estação que define o cenário de implantação modesta tem a seguinte configuração (especificação completa no Apêndice~\ref{apx:specs}):

\begin{itemize}
  \item \textbf{CPU}: Intel Core i7-12700H (12.ª geração, 14 núcleos / 20 \textit{threads}, base 2,3~GHz);
  \item \textbf{RAM}: 32~GB DDR4 a 3200~MHz;
  \item \textbf{GPU dedicada}: NVIDIA GeForce RTX 3050 Laptop GPU com 4~GB de VRAM;
  \item \textbf{Disco}: SSD NVMe de 512~GB;
  \item \textbf{Sistema operacional}: Microsoft Windows 11 Pro (versão 10.0.26200, 64 bits);
  \item \textbf{Software}: Python 3.14.4; Ollama 0.34.0; LangChain 1.x (\texttt{langchain-core} 1.6.4, \texttt{langchain-ollama} 1.1.0, cliente \texttt{ollama} 0.6.2); PostgreSQL 16.4 com PostGIS 3.4 em contêiner Docker, usado apenas na demonstração ponta a ponta.
\end{itemize}

Com 4~GB de VRAM, nenhum dos modelos avaliados executa integralmente em GPU: mesmo o Qwen 3 4B, o menor deles, teve \res{estacao}{qwen}{nagpu} do modelo carregado (pesos e contexto) alocado na VRAM, e o Gemma 4 E2B, \res{estacao}{e2b}{nagpu}; o restante executa na CPU, com penalidade de latência (Seção~\ref{sec:validacao-local}).

\paragraph{Rodadas completas.} As rodadas completas dos quatro modelos locais, com três repetições por modelo, foram executadas numa GPU de nuvem NVIDIA T4 com 16~GB nominais de VRAM (15~GB utilizáveis, conforme o manifesto e a Tabela~\ref{tab:ambiente}), suficiente para cada modelo inteiro, com Ollama (na versão registrada no manifesto de cada rodada), as \textit{tags} e quantizações da Seção~\ref{sec:selecao-modelos} (para o Qwen 3 4B e o Gemma 4 E2B, as mesmas da validação na estação de referência), o mesmo \textit{prompt} de sistema e a mesma definição da ferramenta, executados pelo comando \texttt{pfc-avaliar}. Isso separa o efeito do modelo do efeito do \textit{offload} e evita ocupar a estação de referência durante horas de execução. A qualidade da tradução depende essencialmente do modelo, e não do \textit{hardware} --- diferenças numéricas entre execução em GPU e em CPU podem alterar respostas isoladas, o que a validação na estação de referência quantifica ---; a latência, ao contrário, depende fortemente do \textit{hardware} e é reportada separadamente por ambiente.

\paragraph{Validação na estação de referência.} As 62 consultas das camadas P e N foram executadas também na estação de referência pelos dois modelos de menor porte (Qwen 3 4B e Gemma 4 E2B), com uma repetição, para medir a latência no cenário de implantação modesta e verificar a concordância das respostas com as rodadas completas. As duas rodadas foram feitas na mesma sessão, primeiro a do Qwen 3 4B; antes da rodada do Gemma 4 E2B, todos os modelos residentes no Ollama foram descarregados. Os dois modelos foram, assim, carregados a frio, com VRAM livre semelhante no início da rodada (\res{estacao}{qwen}{vramlivre} e \res{estacao}{e2b}{vramlivre}). O manifesto registra a VRAM livre no início da rodada e a parcela do modelo que o Ollama alocou na GPU, porque a latência nessa estação depende diretamente desses dois valores.

\paragraph{Referência em nuvem.} Os \res{groq}{geral}{nmodelos} modelos do Groq (Seção~\ref{sec:selecao-modelos}) executaram as \res{groq}{geral}{nexecutadas} consultas com uma repetição, respeitando os limites de taxa do provedor.

\subsection{Configuração dos Modelos}
\label{sec:config-modelos}

Todos os modelos foram avaliados com a mesma configuração: \textbf{temperatura zero}; limite de \textbf{1024 \textit{tokens}} de saída; \textbf{nenhum contexto de conversações anteriores}; \textbf{nenhum exemplo resolvido no \textit{prompt}} (\textit{zero-shot}) e \textbf{nenhum dicionário de normalização}; o mesmo \textit{prompt} de sistema, com a data de referência injetada; e a mesma definição da ferramenta \texttt{buscar\_catalogo} (as especificações v2 e v3, avaliadas só com o Gemma 4 E4B, mudam o \textit{prompt} e as descrições; Seções~\ref{sec:lote-validacao} e~\ref{sec:v3}). Nas configurações v1 e v2, não há nova tentativa nem mecanismo de \textit{fallback}: cada falha do modelo é registrada e contada. A v3 acrescenta, por desenho, o retorno de um validador e, no controle com Saída Estruturada, novas tentativas; cada retorno fica registrado no traço da execução, e a análise separa a primeira decisão da resposta final.

\paragraph{Modo \textit{thinking}.} O modo de raciocínio interno foi desativado por parâmetro (\texttt{think=false} na API do Ollama) nos modelos que o declaram, para comparabilidade das latências. Na \textit{build} padrão do Qwen 3 (\texttt{qwen3:4b}) não foi possível desativá-lo de forma limpa --- nem pela diretiva \texttt{/no\_think} no \textit{prompt} nem pelo parâmetro \texttt{think=false} ---: nos testes preliminares, o modelo gerou centenas de \textit{tokens} de raciocínio por consulta, com latências de dezenas de segundos, e, nas consultas em que esgotou o limite de 1024 \textit{tokens} de saída, não chegou a emitir a chamada de ferramenta. Adotou-se, por isso, a \textit{build} oficial sem raciocínio (\textit{Qwen3-4B-Instruct-2507}), de mesmo porte e mesma quantização. Na referência em nuvem, o esforço de raciocínio foi fixado no mínimo aceito por cada modelo \cite{groq2026reasoning} --- \res{groq}{geral}{esforcos}; o valor \texttt{none} desliga o raciocínio, e a família GPT-OSS não permite desligá-lo ---, o que é registrado no manifesto.

\subsection{Linha de Base: Saída Estruturada}
\label{sec:linhas-de-base}

Para medir o que a troca da Saída Estruturada pelo \textit{Tool Calling} rendeu, os modelos foram avaliados também com a abordagem do protótipo, sobre as mesmas consultas e com as mesmas métricas. Nessa linha de base, a chamada de ferramenta é substituída por um pedido de resposta estruturada: o \textit{prompt} de sistema é o mesmo, sem as duas instruções que tratam de chamar ou não a ferramenta, e passa a pedir um objeto JSON no formato do \textit{schema} dos parâmetros, incluído no próprio \textit{prompt} com as mesmas descrições e os mesmos enumerados da ferramenta; o Ollama é usado no modo JSON. A temperatura, o limite de \textit{tokens}, a data de referência e a ausência de exemplos e de dicionário são os da configuração avaliada. Duas diferenças decorrem da própria abordagem: o modelo não tem como recusar --- toda consulta resulta em busca, como no protótipo ---, e os enumerados não são impostos por uma ferramenta, de modo que uma grafia fora da forma canônica conta como erro. Valores vazios que o modelo emite para campos ausentes (texto vazio, limite zero), e que a aplicação ignoraria, são tratados como ausentes, o que favorece a linha de base. A decodificação restrita ao \textit{schema}, que o Ollama também oferece, foi testada e descartada: com ela, os modelos fechavam os intervalos de datas logo depois do primeiro limite, um artefato da gramática que penalizaria a linha de base.

\paragraph{Método do protótipo.} Uma segunda linha de base reproduz a solução anterior tal como foi entregue, portada para Python a partir do código do protótipo: o modelo Phi-4 14B; o dicionário de normalização aplicado à consulta antes do modelo, com as suas 159 entradas; o \textit{prompt} original, em inglês, com dez exemplos resolvidos; o \textit{schema} com o campo obrigatório de justificativa e os enumerados do protótipo, validado a cada resposta, com até três novas tentativas quando a validação falha; e o \textit{fallback} por expressões regulares. Três diferenças são deliberadas: a temperatura é zero, e não 0,6; os valores padrão de ordenação que o protótipo acrescenta a toda busca não contam como parâmetros extraídos; e a validação de estados e municípios contra o banco não é feita. A saída é traduzida para o vocabulário do PFC pela mesma correspondência usada para alinhar o gabarito da camada P (``SCN Carta Topografica Matricial'' para ``SCN Carta Topográfica Matricial'', por exemplo).

\paragraph{Execução.} As linhas de base foram executadas com uma repetição: as dos modelos locais e o método do protótipo, na mesma T4 das rodadas completas, pelo caderno do Colab do repositório; a Saída Estruturada dos modelos em nuvem, no Groq, com a data de referência e o esforço de raciocínio das rodadas com \textit{Tool Calling} de cada modelo. Os dois modelos de menor porte foram comparados também na estação de referência, sobre as camadas P e N.

\subsection{Métricas}
\label{sec:metricas}

\paragraph{Comparação.} A resposta do modelo é comparada com o gabarito campo a campo (Capítulo~\ref{cap:referencial}): enumerados exigem igualdade exata; \textit{strings} livres (\texttt{keyword}, \texttt{state}, \texttt{city}) são comparadas após normalização de caixa, acentos e pontuação; períodos exigem coincidência exata dos limites; \texttt{limit} é inteiro. Quando o gabarito tem leituras múltiplas, a resposta é comparada com cada leitura aceita e avaliada contra a de melhor casamento: correta se coincidir com alguma; por campo, contra a leitura que maximiza acertos menos erros. Campos opcionais ausentes não contam como erro.

\paragraph{Acurácia por consulta.} Uma consulta é correta se o conjunto de parâmetros emitidos coincide com uma leitura aceita --- nem campos a mais, nem a menos. Nas consultas cujo gabarito é não chamar a ferramenta, correta é a resposta sem \textit{tool call}. A acurácia é a proporção de execuções corretas e é reportada com intervalo de confiança de 95\% pelo método de Wilson~\cite{wilson1927probable}, adequado a proporções com amostras moderadas. Como as repetições de uma mesma consulta não são observações independentes --- numa mesma sessão, elas quase sempre coincidem ---, o intervalo usa como tamanho da amostra o número de consultas, e não o de execuções.

\paragraph{Precisão, \textit{recall} e F1 por campo.} Para cada campo: verdadeiro positivo (TP) quando o campo é emitido e corresponde a um valor aceito; falso positivo (FP) quando é emitido e não corresponde, ou não deveria ter sido emitido; falso negativo (FN) quando é obrigatório no gabarito e não foi emitido. Um valor errado conta, portanto, como FP. O F1 médio é ponderado pelo número de ocorrências de cada campo no gabarito; as médias \textit{macro} e \textit{micro} ficam registradas no relatório de cada rodada, no repositório.

\paragraph{Latência.} Tempo da chamada de tradução: do envio da consulta ao modelo até o recebimento da resposta com os \textit{tool calls}, medido com relógio monotônico. Na v3 e no seu controle, que podem fazer mais de uma chamada, vai do envio da consulta à decisão final e inclui todas as chamadas ao modelo e a execução das ferramentas auxiliares e do validador; na arquitetura em duas etapas simulada, é a soma das duas chamadas quando há busca e só a da decisão quando não há. O tempo de execução da ferramenta (SQL) é registrado à parte e não entra na comparação entre modelos. Reportam-se mínimo, mediana, média, percentil 95 e máximo, por ambiente de execução.

\paragraph{Comparação entre modelos.} Diferenças de acurácia entre dois modelos, sobre as mesmas consultas, são testadas pelo teste de McNemar exato~\cite{mcnemar1947note}, adequado para comparar dois classificadores avaliados sobre as mesmas instâncias~\cite{dietterich1998approximate}; o teste usa apenas os pares discordantes (consultas em que só um dos modelos acerta). Pela mesma razão, o teste usa uma observação por consulta: o resultado de cada consulta é o da maioria das suas repetições.

\paragraph{Análise de erros.} Cada resposta incorreta é classificada em: campo inventado (emitido sem estar no gabarito), campo omitido, valor errado, erro misto, não chamou a ferramenta quando devia, chamou quando não devia e argumentos fora do \textit{schema} (tipo inválido ou campo inexistente). Erros de infraestrutura (\textit{timeout}, indisponibilidade) são separados dos erros do modelo e não entram nas estatísticas de latência. Como discutido no Capítulo~\ref{cap:referencial}, \textbf{não se busca um ranking único}: cada métrica revela um aspecto da tradução, e a recomendação final pondera qualidade e latência.

\subsection{Reprodutibilidade}
\label{sec:reprodutibilidade}

Cada rodada registra em manifesto: \textit{tag} e \textit{digest} do modelo no Ollama, versões do Ollama e do LangChain, configuração de inferência, data de referência, \textit{hardware} e os \textit{hashes} SHA-256 do \textit{dataset} e da definição da ferramenta. A pontuação é sempre recalculada contra o gabarito vigente: uma correção de gabarito não exige nova execução dos modelos, e uma consulta cujo texto tenha mudado é descartada da pontuação, com registro. O \textit{dataset}, o gerador, o manual, os relatórios de auditoria, os registros de execução e os \textit{scripts} estão no repositório do projeto (Apêndice~\ref{apendice:codigo}).
